\documentclass[10pt,a4paper]{amsart}
\usepackage[margin=19mm]{geometry}
\usepackage{amsmath,amssymb,mathtools}
\usepackage[scr=rsfs,cal=euler]{mathalfa}
\usepackage{xfrac}
\usepackage[unicode,hidelinks,bookmarks=false]{hyperref}
\newcommand{\ie}{i.e.\ }
\newcommand{\cf}{cf.\ }
\newcommand{\resp}{respectively\ }
\newcommand{\Subsection}{\subsection}
\providecommand{\nobreakdash}{\nobreak\hbox{-}}
\setlength{\emergencystretch}{2em}
\multlinegap0pt
\usepackage{amssymb}
\usepackage{tikz}
\newtheorem{lemma}{Lemma}
\title{Introduction to Dieudonn\'e modules and supersingular abelian varieties revisited}
\author{Chia-Fu Yu}
\date{Source: arXiv:2603.11506v1 (CC BY 4.0)}
\begin{document}
\maketitle

\noindent\emph{Source and licence.} Introduction to Dieudonn\'e modules and supersingular abelian varieties revisited. Version arXiv:2603.11506v1 (CC BY 4.0), distributed by arXiv under the Creative Commons Attribution 4.0 International licence, http://creativecommons.org/licenses/by/4.0/, as recorded in the arXiv licence metadata for this exact version and shown on the abstract page https://arxiv.org/abs/2603.11506v1. Full licence notice: \texttt{licenses/arxiv-2603.11506-CC-BY-4.0.txt}.\par\smallskip

\noindent\textit{Editorial scope.} Counted unit: the article's lemma lm:4.2 (source0.tex lines 751--757). The article's complete original proof of it is delivered with it (source0.tex lines 758--792, one \texttt{proof} environment of 35 lines). No mathematics is written, repaired or re-derived: the statement and the proof are the article's own lines, in the article's own notation, and no retained line is substituted or re-flowed.  No further source-local context is needed: every cross-reference of every retained line resolves inside the two delivered spans.  Nothing was omitted from any retained span, so the package carries no omission record.  No retained line contains a citation, so the wrapper carries no bibliography and no bibliography entry.  No retained line contains a figure environment, an image-inclusion command or drawing code, so no image is delivered and the package declares no asset dependency.  Editorial formatting only, in addition: an AMS \texttt{amsart} preamble that restates the article's own declarations and macros used by the retained lines, namely the environments \texttt{lemma} on separate counters, together with the standard packages those lines need.  The retained environments print in this document as Lemma 1 in order of appearance, whereas the article prints the same items with the numbering of its own full document.  The complete native source of the article as distributed in the arXiv e-print for this exact version is bundled unchanged as \texttt{sources/196221/source0.tex}.
\begin{lemma} \label{lm:4.2}
   Let $F^n+a_1 F^{n-1}+\dots +a_n\in W[F]$ with $n\ge 1$. Then there exist co-prime integers $r\ge 1$ and $s$, and elements $b_0,\dots ,b_{n-1}, u\in W[p^{1/r}]$ with $u$ invertible such that 
   \begin{equation}\label{eq:4.2}
       F^n+a_1 F^{n-1}+\dots +a_n=(b_0 F^{n-1}+\dots +b_{n-1})(F-p^{s/r}) u 
   \end{equation}
   in $W[p^{1/r}][F]$.
\end{lemma}

\begin{proof}
    Let $\lambda:=\min \{\frac{v(a_1)}{1},\dots \frac{v(a_n)}{n} \}$. Then $i\cdot \lambda\le v(a_i)$ for all integers $1\le i \le n$, and $i\cdot \lambda=v(a_i)$ for some $i$. Write $\lambda=s/r$ and $a_i= p^{is/r} \alpha_i$ for each $i$, then $\alpha_i\in W[p^{1/r}]$ and $\alpha_i$ is a unit for some $i$. Put $b_i=p^{is/r} \beta_i$, $1\le i \le n$, and $v=u^{-1}$. Multiplying $v$ by the left hand side of \eqref{eq:4.2}, one has
    \[ (F^n+a_1 F^{n-1}+\dots +a_n)v=v^{(n)} F^n+v^{(n-1)} a_1 F^{n-1} +\dots+v a_n. \]
    The corresponding right hand side is  
    \[ b_0 F^n+(b_1-p^{s/r} b_0)F^{n-1}+(b_2-p^{s/r} b_1) F^{n-2}+ \dots + (b_{n-1}-p^{s/r} b_{n-2})F-b_{n-1} p^{s/r}. \]
    Comparing the coefficients, one obtains
    \[ v^{(n)}=b_0, \quad v^{(n-1)} a_1=b_1-p^{s/r} b_0,\dots, \]
    \[ v^{(1)} a_{n-1}=b_{n-1}-p^{s/r} b_{n-2}, \quad va_n=-p^{s/r} b_{n-1},\]
    or equivalently,
    \begin{equation} \label{eq:4.21}   
     \begin{split}
        v^{(n)} &=\beta_0,\\
        \alpha_1 v^{(n-1)}& =\beta_1-\beta_0,\\
           & \vdots \\
           \alpha_{n-1} v^{(1)}&=\beta_{n-1}-\beta_{n-2}\\
           \alpha_n v&=-\beta_{n-1}.
     \end{split}
     \end{equation}
Thus, Equations~\eqref{eq:4.21} are solvable if and only there exists a unit $v\in W[p^{1/r}]$ such that 
\begin{equation}\label{eq:4.22}
   P(v):=v^{(n)}+ \alpha_1 v^{(n-1)} +\dots+\alpha_{n-1}v^{(1)}+\alpha_n v=0. 
\end{equation}
We solve Equation~\eqref{eq:4.22} by successive approximation. Clearly, the equation 
$P(\bar v)=\bar v^{p^n}+\bar \alpha_1 \bar v^{p^{n-1}}+\dots +\bar \alpha_n \bar v=0$
%where $\bar v$ and $\bar \alpha_i$ are the respective images of $v$ and $\alpha_i$ in $k$, 
has a nonzero solution $\bar v$ in $k$ as $\bar \alpha_i\neq 0$ for some $i$. Thus, one has $P(v_1)\equiv 0 \mod p^{1/r}$, where $v_1$ is a lift of $\bar v$.


Suppose we have $P(v_k)=z p^{k/r}$ for some unit $v_k\in W[p^{1/r}]$.
Put $v_{k+1}=v_k+p^{k/r}x$, where $x\in W[p^{1/r}]$. We need to solve
\[ P(v_{k+1})=p^{k/r}(P(x)+z) \equiv 0 \mod p^{(k+1)/r}, \]
or equivalently to solve the equation 
\[ \bar x^{p^n}+\bar \alpha_1 \bar x^{p^{n-1}}+\dots +\bar \alpha_n \bar x +\bar z=0\]
in $k$, which is solvable. Since $W[p^{1/r}]$ is complete, the equation $P(v)=0$ has a unit solution. \qed
\end{proof}

\end{document}
